\documentclass[10pt,a4paper]{amsart}
\usepackage[margin=19mm]{geometry}
\usepackage{amsmath,amssymb,mathtools}
\usepackage[scr=rsfs,cal=euler]{mathalfa}
\usepackage{xfrac}
\usepackage[unicode,hidelinks,bookmarks=false]{hyperref}
\newcommand{\ie}{i.e.\ }
\newcommand{\cf}{cf.\ }
\newcommand{\resp}{respectively\ }
\newcommand{\Subsection}{\subsection}
\providecommand{\nobreakdash}{\nobreak\hbox{-}}
\setlength{\emergencystretch}{2em}
\multlinegap0pt
\usepackage{amssymb}
\usepackage{mathtools}
\usepackage{booktabs}
\usepackage{tikz}
\usepackage[shortlabels]{enumitem}
\newtheorem{theorem}{Theorem}
\title{Four Dominion Growth Regimes in Trees: Forcing, Fibonacci Enumeration, Periodicity, and Stability}
\author{Julian Allagan, Erin Gray, Jennifer Sawyer,, Gabrielle Morgan}
\date{Source: arXiv:2601.03485v1 (CC BY 4.0)}
\begin{document}
\maketitle

\noindent\emph{Source and licence.} Four Dominion Growth Regimes in Trees: Forcing, Fibonacci Enumeration, Periodicity, and Stability. Version arXiv:2601.03485v1 (CC BY 4.0), distributed by arXiv under the Creative Commons Attribution 4.0 International licence, http://creativecommons.org/licenses/by/4.0/, as recorded in the arXiv licence metadata for this exact version and shown on the abstract page https://arxiv.org/abs/2601.03485v1. Full licence notice: \texttt{licenses/arxiv-2601.03485-CC-BY-4.0.txt}.\par\smallskip

\noindent\textit{Editorial scope.} Counted unit: the article's theorem thm:alt (source0.tex lines 277--298). The article's complete original proof of it is delivered with it (source0.tex lines 300--385, one \texttt{proof} environment of 86 lines). No mathematics is written, repaired or re-derived: the statement and the proof are the article's own lines, in the article's own notation, and no retained line is substituted or re-flowed.  No further source-local context is needed: every cross-reference of every retained line resolves inside the two delivered spans.  Nothing was omitted from any retained span, so the package carries no omission record.  No retained line contains a citation, so the wrapper carries no bibliography and no bibliography entry.  No retained line contains a figure environment, an image-inclusion command or drawing code, so no image is delivered and the package declares no asset dependency.  Editorial formatting only, in addition: an AMS \texttt{amsart} preamble that restates the article's own declarations and macros used by the retained lines, namely the environments \texttt{theorem} on separate counters, together with the standard packages those lines need.  The retained environments print in this document as Theorem 1 in order of appearance, whereas the article prints the same items with the numbering of its own full document.  The complete native source of the article as distributed in the arXiv e-print for this exact version is bundled unchanged as \texttt{sources/192228/source0.tex}.
\begin{theorem}[Fibonacci dominion for alternating combs]\label{thm:alt}
Let $n\ge 2$ and set $k=\lfloor n/2\rfloor$. Then
\[
\gamma(E_n)=k,\qquad
\zeta(E_n)=
\begin{cases}
F_{k+1},& n=2k,\\
F_k,& n=2k+1,
\end{cases}
\]
and
\[
\gamma(O_n)=\lceil n/2\rceil,\qquad
\zeta(O_n)=
\begin{cases}
F_{k+1},& n=2k,\\
F_{k+3},& n=2k+1,
\end{cases}
\]
where $(F_t)_{t\ge1}$ denotes the Fibonacci sequence defined by $F_1=F_2=1$ and
$F_{t+1}=F_t+F_{t-1}$.
\end{theorem}

\begin{proof}
We treat the even and odd attachment patterns separately.

\smallskip
\noindent\emph{Domination numbers.}
In $E_n$, each even index $2i$ contributes a pendant cluster
$C_{2i}=\{v_{2i},\ell_{2i}\}$ with $N[\ell_{2i}]=C_{2i}$.
Every dominating set must intersect each $C_{2i}$, and the clusters are pairwise disjoint.
Hence $\gamma(E_n)\ge k=\lfloor n/2\rfloor$.
Selecting all even hosts $\{v_2,v_4,\dots\}$ dominates $E_n$, so $\gamma(E_n)=k$.
The same argument, applied to odd indices, yields $\gamma(O_n)=\lceil n/2\rceil$.

\smallskip
\noindent\emph{Even attachment.}
For $t\ge 1$, let $H_t$ be the induced subgraph of $E_{2t}$ on path vertices
$v_1,\dots,v_{2t}$ together with pendants $\ell_2,\ell_4,\dots,\ell_{2t}$,
and write $a_t=\zeta(H_t)$.
Since $H_t$ contains exactly $t$ disjoint clusters $C_{2i}=\{v_{2i},\ell_{2i}\}$ and
$\gamma(H_t)=t$, every minimum dominating set intersects each cluster in exactly one vertex.

The initial values follow by inspection.
For $t=1$, $H_1$ is the path $v_1v_2$ with a pendant at $v_2$; the unique minimum
dominating set is $\{v_2\}$, so $a_1=1$.
For $t=2$, $H_2$ is the path $v_1v_2v_3v_4$ with pendants at $v_2$ and $v_4$.
Exactly two minimum dominating sets exist, namely $\{v_2,v_4\}$ and $\{v_2,\ell_4\}$,
so $a_2=2$.

Now fix $t\ge 3$ and let $D$ be a minimum dominating set of $H_t$.
Consider the final cluster $C_{2t}=\{v_{2t},\ell_{2t}\}$.

\emph{Case~1: $v_{2t}\in D$.}
Removing the vertices $\{v_{2t-1},v_{2t},\ell_{2t}\}$ and restricting $D$ to the
remaining induced subgraph yields a minimum dominating set of $H_{t-1}$.
Conversely, any minimum dominating set of $H_{t-1}$ extends uniquely to one of $H_t$
by adding $v_{2t}$.
Thus this case contributes $a_{t-1}$ choices.

\emph{Case~2: $\ell_{2t}\in D$.}
Then $v_{2t}\notin D$, and $v_{2t-1}$ must be dominated by a neighbor in $D$.
Since $v_{2t-1}$ is adjacent only to $v_{2t-2}$ and $v_{2t}$, it follows that
$v_{2t-2}\in D$ is forced.
Removing the last two even clusters and restricting to the induced subgraph on
$v_1,\dots,v_{2t-4}$ yields a minimum dominating set of $H_{t-2}$, and this correspondence
is bijective.
Hence this case contributes $a_{t-2}$ choices.

Combining the two cases gives the recurrence $a_t=a_{t-1}+a_{t-2}$ for $t\ge 3$,
with $a_1=1$ and $a_2=2$, and therefore $a_t=F_{t+1}$.
If $n=2k$, then $E_n=H_k$ and $\zeta(E_{2k})=F_{k+1}$.
If $n=2k+1$, the endpoint $v_{2k+1}$ has no pendant and forces $v_{2k}$ in every
minimum dominating set, reducing to $H_{k-1}$ and yielding $\zeta(E_{2k+1})=F_k$.

\smallskip
\noindent\emph{Odd attachment.}
For $t\ge 0$, let $J_t$ be the induced subgraph of $O_{2t+1}$ on vertices
$v_1,\dots,v_{2t+1}$ together with pendants $\ell_1,\ell_3,\dots,\ell_{2t+1}$,
and write $b_t=\zeta(J_t)$.
The graph $J_t$ contains $t+1$ disjoint clusters
$\{v_{2i-1},\ell_{2i-1}\}$, so $\gamma(J_t)=t+1$.

For $t=0$, $J_0$ consists of $v_1$ with a pendant $\ell_1$, and the unique minimum
dominating set is $\{v_1\}$; thus $b_0=1$.
For $t=1$, $J_1$ is the path $v_1v_2v_3$ with pendants at $v_1$ and $v_3$.
Exactly three minimum dominating sets exist, namely
$\{v_1,v_3\}$, $\{v_1,\ell_3\}$, and $\{\ell_1,v_3\}$, so $b_1=3$.

Fix $t\ge 2$ and let $D$ be a minimum dominating set of $J_t$.
Consider the final cluster $\{v_{2t+1},\ell_{2t+1}\}$.

\emph{Case~1: $v_{2t+1}\in D$.}
Removing $\{v_{2t},v_{2t+1},\ell_{2t+1}\}$ and restricting $D$ yields a minimum
dominating set of $J_{t-1}$, and this operation is reversible.
Hence this case contributes $b_{t-1}$ choices.

\emph{Case~2: $\ell_{2t+1}\in D$.}
Then $v_{2t+1}\notin D$, and domination of $v_{2t}$ forces $v_{2t-1}\in D$.
After fixing these vertices, the remaining choices lie in the induced subgraph
isomorphic to $J_{t-2}$, again bijectively.
Thus this case contributes $b_{t-2}$ choices.

Consequently $b_t=b_{t-1}+b_{t-2}$ for $t\ge 2$, with $b_0=1$ and $b_1=3$,
so $b_t=F_{t+3}$.
If $n=2k+1$, then $O_n=J_k$ and $\zeta(O_{2k+1})=F_{k+3}$.
If $n=2k$, the unpended endpoint $v_{2k}$ forces $v_{2k-1}$ in every minimum
dominating set, reducing to $J_{k-1}$ and giving $\zeta(O_{2k})=F_{k+1}$.
\end{proof}

\end{document}
